
\subsection{由哈密顿同痕移动拉格朗日子流形}
\label{subsec:5-3-2}

本小节考察分别用两个不同的哈密顿同痕
移动两条拉格朗日子流形的情形。
设
$\psi^{(i)}=\{\psi^{(i)}_\rho\}_\rho$
（$\rho\in[0,1]$，$i=0,1$）
为满足 $\psi^{(i)}_0=\mathrm{id}$
的哈密顿同痕。
令 $(L^{(0)},L^{(1)})$
是一对相对 spin 的拉格朗日子流形，并置
\begin{equation}
L^{(i)\prime}=\psi^{(i)}_1(L^{(i)}).
\tag{5.3.13}\label{eq:5-3-13}
\end{equation}
在 (5.3.13) 中，
$L^{(0)}$、$L^{(1)}$
是由不同的辛微分同胚
$\psi^{(0)}_1$、$\psi^{(1)}_1$
移动的；
这一点使本小节的论证
不同于上一小节。

我们允许两条拉格朗日子流形
不横截，而只要求干净相交。
具体地，假设
$L^{(0)}\cap L^{(1)}$ 与
$L^{(0)\prime}\cap L^{(1)\prime}$
均为干净相交，并置
\[
L^{(0)}\cap L^{(1)}
=\coprod_{h\in\pi_0(L^{(0)}\cap L^{(1)})}R_h,
\qquad
L^{(0)\prime}\cap L^{(1)\prime}
=\coprod_{h'\in\pi_0(L^{(0)\prime}\cap L^{(1)\prime})}R'_{h'}.
\]
令 $\Theta^-_{R_h}$、
$\Theta^-_{R'_{h'}}$
为第 3.7.5 小节引入的实线丛。

取相容殆复结构族
$\{J_t\}_t$、$\{J'_t\}_t$
（$t\in[0,1]$）。
由 $J_0,J'_0,J_1,J'_1$，
链的可数集
$\mathcal X_1(L^{(i)})$、
$\mathcal X_2(L^{(i)\prime})$
以及适当模空间上的多值截面
$\mathfrak s_i,\mathfrak s'_i$，
得到滤过 $A_\infty$ 代数
\[
(C_1(L^{(i)};\Lambda_{0,\mathrm{nov}}),
 \mathfrak m_k^{J_i,\mathfrak s_i}),
\qquad
(C_2(L^{(i)\prime};\Lambda_{0,\mathrm{nov}}),
 \mathfrak m_k^{J'_i,\mathfrak s'_i}).
\]
按照第 3.7.5 小节，
取以 $\Theta^-_{R_h}$、
$\Theta^-_{R'_{h'}}$ 为系数的
光滑奇异单形的可数集
$\mathcal X_1(R_h)$、
$\mathcal X_2(R'_{h'})$。
由定理 3.7.72，得到滤过
$A_\infty$ 双模
\[
\begin{aligned}
&(C_1(L^{(1)},L^{(0)};
 \Lambda_{0,\mathrm{nov}}),
 \mathfrak n_{k_1,k_0}^{J,\mathfrak s}),
\\
&(C_2(L^{(1)\prime},L^{(0)\prime};
 \Lambda_{0,\mathrm{nov}}),
 \mathfrak n_{k_1,k_0}^{J',\mathfrak s'}).
\end{aligned}
\]
稍后还要取合适的可数集
$\mathcal X_3(R'_{h'})
\supset\mathcal X_2(R'_{h'})$。
那时由 $\mathcal X_3(R'_{h'})$
得到的滤过 $A_\infty$ 双模记作
\[
(C_3(L^{(1)\prime},L^{(0)\prime};
 \Lambda_{0,\mathrm{nov}}),
 \mathfrak n_{k_1,k_0}^{J',\mathfrak s'}).
\]

取两族相容殆复结构的同伦
$\{J^{(i)}_\rho\}_\rho$，
分别从 $J_i^{\psi^{(i)}_1}$
连接到 $J'_i$（$i=0,1$）：
\begin{center}
\[
\begin{array}{ccc}
J_1^{\psi^{(1)}_1}=J^{(1)}_0
 &\xrightarrow{\;J^{(1)}_\rho\;}&
J^{(1)}_1=J'_1\\[2mm]
\Big\updownarrow\vcenter{\hbox{\scriptsize$J_t$}}
 &&\Big\updownarrow\vcenter{\hbox{\scriptsize$J'_t$}}\\[2mm]
J_0^{\psi^{(0)}_1}=J^{(0)}_0
 &\xrightarrow{\;J^{(0)}_\rho\;}&
J^{(0)}_1=J'_0 .
\end{array}
\]
\renewcommand{\thefigure}{5.3.2}
\refstepcounter{figure}
{\normalfont\bfseries 图式~\thefigure}\label{diag:5-3-2}
\end{center}
连同记为
$\mathfrak s_i^{\operatorname{top}(\rho)}$
的多值截面的选择，
按照第 4.2.2 小节得到滤过
$A_\infty$ 同态
\[
\begin{aligned}
\mathfrak g^{i,\psi^{(i)}_1,\{J^{(i)}_\rho\}_\rho,
 \mathfrak s_i^{\operatorname{top}(\rho)}}:
&(C_1(L^{(i)};\Lambda_{0,\mathrm{nov}}),
 \mathfrak m_k^{J_i,\mathfrak s_i})
\\
&\longrightarrow
(C_2(L^{(i)\prime};\Lambda_{0,\mathrm{nov}}),
 \mathfrak m_k^{J'_i,\mathfrak s'_i}).
\end{aligned}
\]
回忆存在
$\mathcal X_3(L^{(i)\prime})
\supset\mathcal X_2(L^{(i)\prime})$
以及滤过 $A_\infty$ 同态
\[
\begin{aligned}
\mathfrak f^{i,\psi^{(i)}_1,\{J^{(i)}_\rho\}_\rho,
 \mathfrak s_i^{\operatorname{top}(\rho)}}:
&(C_1(L^{(i)};\Lambda_{0,\mathrm{nov}}),
 \mathfrak m_k^{J_i,\mathfrak s_i})
\\
&\longrightarrow
(C_3(L^{(i)\prime};\Lambda_{0,\mathrm{nov}}),
 \mathfrak m_k^{J'_i,\mathfrak s'_i}).
\end{aligned}
\]
前述 $\mathfrak g$ 是此 $\mathfrak f$
与由包含
$\mathcal X_2(L^{(i)\prime})
\subset\mathcal X_3(L^{(i)\prime})$
诱导的滤过 $A_\infty$ 同态之
同伦逆的复合。

用
$\|\{\psi^{(i)}_\rho\}_\rho\|$
表示从恒等映射出发的
哈密顿同痕
$\{\psi^{(i)}_\rho\}_\rho$
的霍费尔长度
（其定义将在第 5.3.5 小节回顾）。
取任意满足
\[
\mathfrak E\ge
\|\{\psi^{(1)}_\rho\}_\rho\|
+\|\{\psi^{(0)}_\rho\}_\rho\|
\]
的数 $\mathfrak E$。
下述定理是第 5.3.2--5.3.5
小节的主要结果。
\begin{theorem}{5.3.14}\label{thm:5-3-14}
存在 $\mathfrak E$-弱滤过
$A_\infty$ 双模同态
\begin{equation}
\begin{aligned}
\varphi\colon
&(C_1(L^{(1)},L^{(0)};\Lambda_{\mathrm{nov}}),
 \mathfrak n_{k_1,k_0}^{J_t,\mathfrak s})
\\
&\longrightarrow
(C_2(L^{(1)\prime},L^{(0)\prime};
 \Lambda_{\mathrm{nov}}),
 \mathfrak n_{k_1,k_0}^{J'_t,\mathfrak s'})
\end{aligned}
\tag{5.3.15}\label{eq:5-3-15}
\end{equation}
立于
$\mathfrak g^{1,\psi^{(1)}_1,
 \{J^{(1)}_\rho\}_\rho,
 \mathfrak s_1^{\operatorname{top}(\rho)}}$
及
$\mathfrak g^{0,\psi^{(0)}_1,
 \{J^{(0)}_\rho\}_\rho,
 \mathfrak s_0^{\operatorname{top}(\rho)}}$
之上。
此外，(5.3.15) 是同伦等价。
\end{theorem}

\begin{remark}{5.3.16}\label{rem:5-3-16}
再次强调，这里以
$\Lambda_{\mathrm{nov}}$ 为系数。
当两条拉格朗日子流形由不同的
哈密顿同痕移动时，
以 $\Lambda_{0,\mathrm{nov}}$
为系数的弗洛尔上同调
\emph{并不}保持不变；
参见第 3.7.6 小节的例子。
\end{remark}

\begin{proof}
本小节构造定理 5.3.14 中的
$A_\infty$ 双模同态 $\varphi$。
在接下来的 5.3.3、5.3.4 小节，
我们再证明它确为同伦等价。

利用上面选取的
$\{J^{(i)}_\rho\}_\rho$，
取相容殆复结构双参数族
$\{J_{\tau,t}\}_{\tau,t}$，
使
\begin{align}
J_{\tau,0}
&=\bigl(J^{(0)}_{\chi(\tau)}\bigr)^{
 \psi^{(0)}_{\chi(\tau)}
 \circ(\psi^{(0)}_1)^{-1}},
\qquad
J_{\tau,1}
=\bigl(J^{(1)}_{\chi(\tau)}\bigr)^{
 \psi^{(1)}_{\chi(\tau)}
 \circ(\psi^{(1)}_1)^{-1}},
\tag{5.3.17.1}\label{eq:5-3-17-1}\\
J_{\tau,t}&=J_t\quad(\tau<-A),
\qquad
J_{\tau,t}=J'_t\quad(\tau>A).
\tag{5.3.17.2}\label{eq:5-3-17-2}
\end{align}
这里
$\bigl(J^{(0)}_\rho\bigr)^{
 \psi^{(0)}_\rho\circ(\psi^{(0)}_1)^{-1}}$
表示用微分同胚
$\psi^{(0)}_\rho\circ(\psi^{(0)}_1)^{-1}$
推送 $J^{(0)}_\rho$ 所得的殆复结构；
$\chi\colon\mathbb R\to(0,1)$
是 (5.3.2) 中的微分同胚，
即 $\rho=\chi(\tau)$。
这两端的结构族可画作
\[
\begin{array}{ccccc}
J_1&&
\underset{\sim}{J_{\tau,1}}
&&J'_1\\
\big\updownarrow\vcenter{\hbox{\scriptsize$J_t$}}
&&&&
\big\updownarrow\vcenter{\hbox{\scriptsize$J'_t$}}\\
J_0&&
\underset{\sim}{J_{\tau,0}}
&&J'_0\\[-1mm]
(\tau<-A)&&&&(\tau>A).
\end{array}
\]

回忆 $\widetilde\pi_0
(L^{(0)}\cap L^{(1)})$
是满足 (3.7.49) 的等价类集合，
见定义 3.7.50。
取
$[h,w]\in\widetilde\pi_0
(L^{(0)}\cap L^{(1)})$、
$[h',w']\in\widetilde\pi_0
(L^{(0)\prime}\cap L^{(1)\prime})$。
同前，令 $\Sigma$ 为
可能带有球面气泡的
$\mathbb R\times[0,1]$。
考察满足以下条件的
$(\Sigma;u;\vec\tau^{(1)},\vec\tau^{(0)})$：
\begin{enumerate}[
 label=\textup{(5.3.18.\arabic*)},
 leftmargin=6.3em]
\item\label{cond:5-3-18-1}
$u\colon\Sigma\to M$
在 $\mathbb R\times[0,1]$ 上满足
\[
\frac{\partial u}{\partial\tau}
+J_{\tau,t}\frac{\partial u}{\partial t}=0.
\]
若在 $(\tau,t)$ 处附着球面气泡树，
则该树中的每个球面都是
$J_{\tau,t}$-全纯的；
此外要求稳定性条件。
\item\label{cond:5-3-18-2}
$u(\tau,0)\in
\psi^{(0)}_{\chi(\tau)}(L^{(0)})$，
$u(\tau,1)\in
\psi^{(1)}_{\chi(\tau)}(L^{(1)})$。
\item\label{cond:5-3-18-3}
$\lim_{\tau\to-\infty}u(\tau)\in R_h$，
$\lim_{\tau\to+\infty}u(\tau)\in R'_{h'}$。
\item\label{cond:5-3-18-4}
$w\#\bigl(
(I^\chi_{\psi^{(0)},\psi^{(1)}})^{-1}u
\bigr)
\sim
(I^\chi_{\psi^{(0)},\psi^{(1)}})^{-1}w'$。
下文解释这一记号。
\item\label{cond:5-3-18-5}
$\vec\tau^{(0)}
=(\tau_1^{(0)},\ldots,\tau_{k_0}^{(0)})$，
其中
$-\infty<\tau_1^{(0)}<\cdots<
 \tau_{k_0}^{(0)}<+\infty$；
$\vec\tau^{(1)}
=(\tau_1^{(1)},\ldots,\tau_{k_1}^{(1)})$，
其中
$-\infty<\tau_{k_1}^{(1)}<\cdots<
 \tau_1^{(1)}<+\infty$。
\end{enumerate}
现解释条件 (5.3.18.4)。
与固定拉格朗日边界的边界映射情形不同，
$u$ 满足移动边界条件；
而 $w,w'$ 所满足的边界条件也不同，
两者由 $\psi^{(0)},\psi^{(1)}$ 联系。
第 5.1.2 小节定义了映射
\[
\begin{aligned}
I_{\vec\psi}
=I_{\psi^{(0)},\psi^{(1)}}\colon
\Omega(L^{(0)},L^{(1)})
\longrightarrow
\Omega\bigl(
\psi^{(0)}_1(L^{(0)}),
\psi^{(1)}_1(L^{(1)})\bigr),
\\[-1mm]
\bigl(I_{\psi^{(0)},\psi^{(1)}}(\ell)\bigr)(t)
=\bigl(\psi^{(1)}_t\circ
\psi^{(0)}_{1-t}\bigr)(\ell(t)).
\end{aligned}
\]
再定义
\[
\bigl(I^\chi_{\psi^{(0)},\psi^{(1)}}(u)
\bigr)(\tau,t)
=\bigl(
\psi^{(1)}_{\chi(\tau)t}
\circ\psi^{(0)}_{\chi(\tau)(1-t)}
\bigr)(u(\tau,t)).
\]
此映射把具有固定边界条件
$(L^{(0)},L^{(1)})$ 的 $u$
变为满足移动边界条件
(5.3.18.2) 的映射。
因此，(5.3.18.4) 中等价关系
两边的边界在 $t=0,1$
分别落于 $L^{(0)},L^{(1)}$。
又由定义可知
\[
\bigl(
(I^\chi_{\psi^{(0)},\psi^{(1)}})^{-1}
(w')\bigr)(0,t)=\ell_0(t),
\]
因为
$w'(0,t)=
\ell_0^{\psi^{(0)},\psi^{(1)}}(t)
=\bigl(
I_{\psi^{(0)},\psi^{(1)}}(\ell_0)
\bigr)(t)$
（参见第 2.2 节）。
另一方面，
\[
\begin{aligned}
\bigl(
(I^\chi_{\psi^{(0)},\psi^{(1)}})^{-1}
(w')\bigr)(1,t)
&=\bigl(
I_{\psi^{(0)},\psi^{(1)}}^{-1}(\ell_q)
\bigr)(t)
\\
&=\bigl(
I_{\psi^{(0)},\psi^{(1)}}^{-1}
(u(\infty,\cdot))
\bigr)(t),
\end{aligned}
\]
这是由 $u$ 在 $\tau=+\infty$
处的渐近条件得到的。
故 (5.3.18.4) 要求的等价关系
确有意义。

记 (5.3.18) 的解所组成的模空间为
\begin{equation}
\begin{aligned}
\mathcal M^\circ_{k_1,k_0}\bigl(
 &(L^{(1)},\psi^{(1)}),
 (L^{(0)},\psi^{(0)});
 \{J_{\tau,t}\}_{\tau,t};
\\[-1mm]
 &[h,w],[h',w'];\operatorname{top}(\vec\tau)
\bigr).
\end{aligned}
\tag{5.3.19}\label{eq:5-3-19}
\end{equation}
要由 (5.3.19) 的模空间
得到弱滤过 $A_\infty$ 双模同态，
还需下面的估计 (5.3.21)。
与边界拉格朗日子流形固定的情形不同，
这里的估计不是自动成立的。
第二署名作者曾在 \textbf{[Oh93]}
作过类似但尚非最优的估计。

\begin{proposition}{5.3.20}\label{prop:5-3-20}
若 (5.3.19) 的模空间非空，则
\begin{equation}
0\le
\mathcal A_{\ell_0^{\vec\psi}}
([\ell_q,w'])
-\mathcal A_{\ell_0}([\ell_p,w])
+\mathfrak E,
\tag{5.3.21.1}\label{eq:5-3-21-1}
\end{equation}
其中 $p\in R_h$、$q\in R'_{h'}$。
特别地，
\begin{equation}
\mathcal A_{\ell_0^{\vec\psi}}
([\ell_q,w'])
-\mathcal A_{\ell_0}([\ell_p,w])
\ge-\mathfrak E.
\tag{5.3.21.2}\label{eq:5-3-21-2}
\end{equation}
此处可取
$\mathfrak E=
\|\psi^{(0)}\|+\|\psi^{(1)}\|$。
\end{proposition}
命题 5.3.20 的证明
将在第 5.3.5 小节给出。

\begin{remark}{5.3.22}\label{rem:5-3-22}
此前各小节研究
固定拉格朗日边界条件下的
柯西--黎曼方程；
在那里，$\mathfrak E=0$
的类似结论相当直接。
我们也曾用它证明所得映射
不仅弱滤过，而且滤过。
\end{remark}

接下来，与上一小节类似，
将 (5.3.19) 的模空间
同第 4.6 节用于定义
$\mathfrak f^{i,\psi^{(i)}_1,
\{J^{(i)}_\rho\}_\rho,
\mathfrak s_i^{\operatorname{top}(\rho)}}$
的模空间取纤维积。
先定义评价映射
\[
\begin{aligned}
\operatorname{ev}
={}&(\operatorname{ev}^{(1)}_1,\ldots,
\operatorname{ev}^{(1)}_{k_1};
\operatorname{ev}_{-\infty};
\operatorname{ev}^{(0)}_1,\ldots,
\operatorname{ev}^{(0)}_{k_0};
\operatorname{ev}_{+\infty})
\\
\colon&
\mathcal M^\circ_{k_1,k_0}
\bigl(
(L^{(1)},\psi^{(1)}),
(L^{(0)},\psi^{(0)});
\{J_{\tau,t}\}_{\tau,t};
[h,w],[h',w'];\operatorname{top}(\vec\tau)
\bigr)
\\
&\longrightarrow
(L^{(1)})^{k_1}\times R_h
\times(L^{(0)})^{k_0}\times R'_{h'}.
\end{aligned}
\]
对系统的评价具体为\translatornote{原书实体页 319
在此展开式中把时间向量次序写成
$(\vec\tau^{(0)},\vec\tau^{(1)})$，
但模空间 (5.3.19) 与下文的系统
$\mathfrak D$ 均按 $(1),(0)$ 排列；
同页时间评价映射的末分量还写成
$k,\ell$，而靶维度分别是
$k_0,k_1$。
译文统一依所定义的系统次序和靶维度排印；
见疑误 A-0300。}
\[
\begin{aligned}
\operatorname{ev}
(\Sigma;u;\vec\tau^{(1)},\vec\tau^{(0)})
=\bigl(
 &(\psi^{(1)}_{\chi(\tau_1^{(1)})})^{-1}
 u(\tau_1^{(1)},1),\ldots,
\\
 &(\psi^{(1)}_{\chi(\tau_{k_1}^{(1)})})^{-1}
 u(\tau_{k_1}^{(1)},1),
 \lim_{\tau\to-\infty}u(\tau),
\\
 &(\psi^{(0)}_{\chi(\tau_1^{(0)})})^{-1}
 u(\tau_1^{(0)},0),\ldots,
\\
 &(\psi^{(0)}_{\chi(\tau_{k_0}^{(0)})})^{-1}
 u(\tau_{k_0}^{(0)},0),
 \lim_{\tau\to+\infty}u(\tau)
 \bigr).
\end{aligned}
\]
还定义时间评价映射
\[
\operatorname{ev}_{\mathrm{time},(0)}
=(\operatorname{ev}_{\mathrm{time},(0),1},
\ldots,\operatorname{ev}_{\mathrm{time},(0),k_0}).
\]
\[
\begin{aligned}
\operatorname{ev}_{\mathrm{time},(0)}\colon&
\mathcal M^\circ_{k_1,k_0}
\bigl((L^{(1)},\psi^{(1)}),
(L^{(0)},\psi^{(0)});
\\[-1mm]
&\{J_{\tau,t}\}_{\tau,t};
[h,w],[h',w'];\operatorname{top}(\vec\tau)\bigr)
\longrightarrow\mathbb R^{k_0},
\\
\operatorname{ev}_{\mathrm{time},(0),j}
(\Sigma;u;\vec\tau^{(1)},\vec\tau^{(0)})
&=\tau_j^{(0)}.
\end{aligned}
\]
同样定义
\[
\operatorname{ev}_{\mathrm{time},(1)}
=(\operatorname{ev}_{\mathrm{time},(1),1},
\ldots,\operatorname{ev}_{\mathrm{time},(1),k_1}).
\]
\[
\begin{aligned}
\operatorname{ev}_{\mathrm{time},(1)}\colon&
\mathcal M^\circ_{k_1,k_0}
\bigl((L^{(1)},\psi^{(1)}),
(L^{(0)},\psi^{(0)});
\\[-1mm]
&\{J_{\tau,t}\}_{\tau,t};
[h,w],[h',w'];\operatorname{top}(\vec\tau)\bigr)
\longrightarrow\mathbb R^{k_1}.
\end{aligned}
\]

取非负整数
$l_1^{(0)},\ldots,l_a^{(0)}$
与 $l_1^{(1)},\ldots,l_b^{(1)}$，
使
\[
\sum_i l_i^{(0)}=k_0,
\qquad
\sum_i l_i^{(1)}=k_1.
\]
考虑系统
\[
\begin{aligned}
\mathfrak D=\bigl(
 &(\Sigma;u;\vec\tau^{(1)},\vec\tau^{(0)});
\\[-1mm]
 &((\Sigma_{(1),i},\vec z_{(1),i}),
   (u_{(1),i,\alpha}),
   (\rho_{(1),i,\alpha}))_{i=1}^{b};
\\[-1mm]
 &((\Sigma_{(0),i},\vec z_{(0),i}),
   (u_{(0),i,\alpha}),
   (\rho_{(0),i,\alpha}))_{i=1}^{a}
\bigr).
\end{aligned}
\]
它须满足下列类似于 (5.3.5) 的条件：
\begin{enumerate}[
 label=\textup{(5.3.23.\arabic*)},
 leftmargin=6.4em]
\item\label{cond:5-3-23-1}
$(\Sigma;u;\vec\tau^{(1)},\vec\tau^{(0)})$
属于
\[
\mathcal M^\circ_{b,a}
\bigl((L^{(1)},\psi^{(1)}),
(L^{(0)},\psi^{(0)});
\{J_{\tau,t}\}_{\tau,t};
[h,w],[h',w'];\operatorname{top}(\vec\tau)\bigr).
\]
\item\label{cond:5-3-23-2}
对下边界每个 $i$，
\[
\begin{aligned}
&((\Sigma_{(0),i},\vec z_{(0),i}),
(u_{(0),i,\alpha}),
(\rho_{(0),i,\alpha}))
\\
&\hspace{2em}\in
\mathcal M^{\mathrm{main}}_{l_i^{(0)}+1}
(M,L^{(0)},\{J^{(0)}_\rho\}_\rho:\beta_i^{(0)};\operatorname{top}(\rho)).
\end{aligned}
\]
\item\label{cond:5-3-23-3}
对上边界每个 $i$，
\[
\begin{aligned}
&((\Sigma_{(1),i},\vec z_{(1),i}),
(u_{(1),i,\alpha}),
(\rho_{(1),i,\alpha}))
\\
&\hspace{2em}\in
\mathcal M^{\mathrm{main}}_{l_i^{(1)}+1}
(M,L^{(1)},\{J^{(1)}_\rho\}_\rho:\beta_i^{(1)};\operatorname{top}(\rho)).
\end{aligned}
\]
\item\label{cond:5-3-23-4}
\[
\begin{aligned}
\operatorname{ev}_0
\bigl((\Sigma_{(0),i},\vec z_{(0),i}),
\\
\qquad(u_{(0),i,\alpha}),
(\rho_{(0),i,\alpha})\bigr)
&=
\bigl(
\psi_1^{(0)}
\circ(\psi_{\chi(\tau_i^{(0)})}^{(0)})^{-1}
\circ u
\bigr)(\tau_i^{(0)},0).
\end{aligned}
\]
\item\label{cond:5-3-23-5}
\[
\begin{aligned}
\operatorname{ev}_0
\bigl((\Sigma_{(1),i},\vec z_{(1),i}),
\\
\qquad(u_{(1),i,\alpha}),
(\rho_{(1),i,\alpha})\bigr)
&=
\bigl(
\psi_1^{(1)}
\circ(\psi_{\chi(\tau_i^{(1)})}^{(1)})^{-1}
\circ u
\bigr)(\tau_i^{(1)},1).
\end{aligned}
\]
\item\label{cond:5-3-23-6}
$\rho_{(0),i,0}\le\chi(\tau_i^{(0)})$，
$\rho_{(1),i,0}\le\chi(\tau_i^{(1)})$。
\item\label{cond:5-3-23-7}
\[
\biggl(\sum_i\beta_i^{(0)}
+\sum_i\beta_i^{(1)}\biggr)
\#w\#
\bigl((I^\chi_{\psi^{(0)},\psi^{(1)}})^{-1}u\bigr)
\sim
(I^\chi_{\psi^{(0)},\psi^{(1)}})^{-1}w'.
\]
\end{enumerate}

回忆 $\rho_{(0),i,0}$
是 $D_{(0),i}$ 最后一个分支
所对应的“时间”\translatornote{原书实体页 320
把此下边界气泡的根分支时间
写成 $\tau_{(1),i,0}$，
下句又写
$J^{(0)}_{\tau_{(0),i,0}}$；
但 (5.3.23.6) 以及上一小节的
同型定义均使用
$\rho_{(0),i,0}$。
这里两处统一按圆盘族参数排印；
见疑误 A-0299。}；
尤其
$u_{(0),i,0}$ 是
$J^{(0)}_{\rho_{(0),i,0}}$-全纯的。
与前面一样，
在记号的宽松用法下还允许
$l_j^{(i)}=1$、
$\beta_j^{(i)}=0$
（$i=0$ 或 $1$）。
此时原先未定义的模空间
\[
\mathcal M^{\mathrm{main}}_2
(M,L^{(i)},\{J^{(i)}_\rho\}_\rho;
0;\operatorname{top}(\rho))
\]
现规定为单点集，
并且 (5.3.23.4)--(5.3.23.6)
在此情形下均无要求。
上述条件除
(5.3.23.4)、(5.3.23.5)
和 (5.3.23.7)
作了适当改动外，
大多与 (5.3.5) 相同。

令
\[
\begin{aligned}
\mathcal M^{\mathrm{unbr}}_{k_1,k_0}
\bigl(
 &(L^{(1)},\psi^{(1)}),
 (L^{(0)},\psi^{(0)});
 \{J_{\tau,t}\}_{\tau,t};
\\[-1mm]
 &[h,w],[h',w'];\operatorname{top}(\vec\tau)
\bigr)
\end{aligned}
\]
为全部此类系统 $\mathfrak D$
所成的集合；
它也可定义为若干模空间纤维积的并。
其紧化定义为
\[
\begin{aligned}
&\mathcal M_{k_1,k_0}
\bigl((L^{(1)},\psi^{(1)}),
(L^{(0)},\psi^{(0)});
\{J_{\tau,t}\}_{\tau,t};
[h,w],[h',w'];\operatorname{top}(\vec\tau)\bigr)
\\
&=
\coprod_{\substack{l_0,l_1\\m_1,m_0}}
\coprod_{\substack{h'',w''\\h''',w'''}}
\biggl[
\mathcal M_{l_1,l_0}
\bigl(\psi^{(1)}_0(L^{(1)}),
\psi^{(0)}_0(L^{(0)});
\{J_{0,t}\}_t;
[h,w],[h'',w'']\bigr)
\\[-1mm]
&\qquad\times_{R_{h''}}
\mathcal M^{\mathrm{unbr}}_{
k_1-l_1-m_1,\,
k_0-l_0-m_0}
\\[-1mm]
&\qquad\quad
\bigl((L^{(1)},\psi^{(1)}),
(L^{(0)},\psi^{(0)});
\{J_{\tau,t}\}_{\tau,t};
[h'',w''],[h''',w'''];
\operatorname{top}(\vec\tau)\bigr)
\\[-1mm]
&\qquad\times_{R_{h'''}}
\mathcal M_{m_1,m_0}
\bigl(\psi^{(1)}_1(L^{(1)}),
\psi^{(0)}_1(L^{(0)});
\{J_{1,t}\}_t;
[h''',w'''],[h',w']\bigr)
\biggr].
\end{aligned}
\]
其中首尾两个模空间
$\mathcal M_{l_1,l_0}$、
$\mathcal M_{m_1,m_0}$
按 (3.7.56) 定义。
与 (5.3.5) 一样，
上述评价映射 $\operatorname{ev}$
可延拓到此紧化：
\[
\begin{aligned}
\operatorname{ev}
=&
(\operatorname{ev}^{(1)}_1,\ldots,
\operatorname{ev}^{(1)}_{k_1};
\operatorname{ev}_{-\infty};
\operatorname{ev}^{(0)}_1,\ldots,
\operatorname{ev}^{(0)}_{k_0};
\operatorname{ev}_{+\infty})
\\
\colon&
\mathcal M_{k_1,k_0}
\bigl((L^{(1)},\psi^{(1)}),
(L^{(0)},\psi^{(0)});
\{J_{\tau,t}\}_{\tau,t};
[h,w],[h',w'];\operatorname{top}(\vec\tau)\bigr)
\\
&\longrightarrow
(L^{(1)})^{k_1}\times R_h
\times(L^{(0)})^{k_0}\times R'_{h'}.
\end{aligned}
\]
现定义模空间
\begin{equation}
\begin{aligned}
\mathcal M_{k_1,k_0}
\bigl(
 &(L^{(1)},\psi^{(1)}),
 (L^{(0)},\psi^{(0)});
 \{J_{\tau,t}\}_{\tau,t};
\\[-1mm]
 &[h,w],[h',w'];\operatorname{top}(\vec\tau);
 \vec P^{(1)},S,\vec P^{(0)}
\bigr)
\end{aligned}
\tag{5.3.24}\label{eq:5-3-24}
\end{equation}
为下列纤维积：
\[
\begin{aligned}
&\mathcal M_{k_1,k_0}
\bigl((L^{(1)},\psi^{(1)}),
(L^{(0)},\psi^{(0)});
\{J_{\tau,t}\}_{\tau,t};
[h,w],[h',w'];\operatorname{top}(\vec\tau)\bigr)
\\
&\qquad{}_{\operatorname{ev}}\!\times
\biggl(
\prod_i P_i^{(1)}
\times S\times\prod_iP_i^{(0)}
\biggr).
\end{aligned}
\]
通过评价映射
$\operatorname{ev}_{+\infty}$，
可将其看作 $R'_{h'}$ 上的一条链。

与命题 5.3.6 类似，
模空间 (5.3.24) 具有 Kuranishi 结构；
可在其上定义与前面已选定的多值截面
相容的多值截面 $\mathfrak s$。
通过适当三角剖分，
其零点集
\[
\begin{aligned}
&\mathcal M_{k_1,k_0}
\bigl((L^{(1)},\psi^{(1)}),
      (L^{(0)},\psi^{(0)});
\{J_{\tau,t}\}_{\tau,t};
[h,w],[h',w'];\\[-1mm]
&\hspace{27mm}
\operatorname{top}(\vec\tau);
\vec P^{(1)},S,\vec P^{(0)}\bigr)^{\mathfrak s}
=\mathfrak s^{-1}(0)
\end{aligned}
\]
给出一条 $\mathbb Q$-链。
定向问题可按第 3.7.5 小节的方法处理。
与上一小节一样，定义
\[
\begin{aligned}
&\phi_{k_1,k_0}
\bigl(T^{\lambda_1^{(1)}}e^{\mu_1^{(1)}}P_1^{(1)},
\ldots,
T^{\lambda_{k_1}^{(1)}}e^{\mu_{k_1}^{(1)}}P_{k_1}^{(1)},
T^\lambda e^\mu S,\\[-1mm]
&\hspace{37mm}
T^{\lambda_1^{(0)}}e^{\mu_1^{(0)}}P_1^{(0)},
\ldots,
T^{\lambda_{k_0}^{(0)}}e^{\mu_{k_0}^{(0)}}P_{k_0}^{(0)}
\bigr)\\
&\quad =
\sum_{\substack{[h',w']\in
\widetilde\pi_0(L^{(0)\prime}\cap L^{(1)\prime})}}
T^{\lambda'}e^{\mu'}
\bigl[
\mathcal M_{k_1,k_0}
\bigl((L^{(1)},\psi^{(1)}),
      (L^{(0)},\psi^{(0)});
\{J_{\tau,t}\}_{\tau,t};\\[-1mm]
&\hspace{39mm}
[h,w],[h',w'];
\operatorname{top}(\vec\tau);
\vec P^{(1)},S,\vec P^{(0)}
\bigr)^{\mathfrak s},
\operatorname{ev}_{+\infty}
\bigr],
\end{aligned}
\]
其中
\[
\lambda'
=\sum_{i=1}^{k_0}\lambda_i^{(0)}
 +\lambda
 +\sum_{i=1}^{k_1}\lambda_i^{(1)},
\qquad
\mu'
=\sum_{i=1}^{k_0}\mu_i^{(0)}
 +\mu
 +\sum_{i=1}^{k_1}\mu_i^{(1)}.
\]

\begin{lemma}{5.3.25}\label{lem:5-3-25}
设 $\phi_{k_1,k_0}$
为上述映射。
\begin{enumerate}[
 label=\textup{(5.3.26.\arabic*)},
 leftmargin=6.3em]
\item\label{cond:5-3-26-1}
可以选取光滑奇异单形的可数集
$\mathcal X_3(R'_{h'})$，
使 $\phi_{k_1,k_0}$ 的像包含于
$C_3(L^{(1)\prime},L^{(0)\prime};
\Lambda_{0,\mathrm{nov}})$。
\item\label{cond:5-3-26-2}
自然包含诱导链同伦等价
\[
C_2(L^{(1)\prime},L^{(0)\prime};\mathbb Q)
\longrightarrow
C_3(L^{(1)\prime},L^{(0)\prime};\mathbb Q).
\]
\item\label{cond:5-3-26-3}
$\phi=\{\phi_{k_1,k_0}\}_{k_1,k_0}$
是立于滤过 $A_\infty$ 同态对
\[
\bigl(
\mathfrak f^{1,\psi^{(1)}_1,
\{J^{(1)}_\rho\}_\rho,
\mathfrak s_1^{\operatorname{top}(\rho)}},
\mathfrak f^{0,\psi^{(0)}_1,
\{J^{(0)}_\rho\}_\rho,
\mathfrak s_0^{\operatorname{top}(\rho)}}
\bigr)
\]
之上的 $\mathfrak E$-弱滤过
$A_\infty$ 双模同态
\[
\phi\colon
C_1(L^{(1)},L^{(0)};\Lambda_{\mathrm{nov}})
\longrightarrow
C_3(L^{(1)\prime},L^{(0)\prime};
\Lambda_{\mathrm{nov}}).
\]
\end{enumerate}
\end{lemma}
\begin{proof}
证明与引理 5.3.8 相同。
注意不等式 (5.3.21.2)
表明 $\phi_{k_1,k_0}$
是 $\mathfrak E$-弱滤过
$A_\infty$ 双模同态。
\end{proof}

由 (5.3.26.2) 及定理 5.2.35，
自然包含
\[
C_2(L^{(1)\prime},L^{(0)\prime};
\Lambda_{0,\mathrm{nov}})
\longrightarrow
C_3(L^{(1)\prime},L^{(0)\prime};
\Lambda_{0,\mathrm{nov}})
\]
是立于包含映射
$C_2(L^{(i)\prime})
\to C_3(L^{(i)\prime})$
之上的同伦等价。
注意这里使用的是
$\Lambda_{0,\mathrm{nov}}$ 系数，
而非 $\Lambda_{\mathrm{nov}}$ 系数；
因此可应用定理 5.2.35，
得到它的同伦逆。
将 $\phi$ 与这个同伦逆复合，
便构造出定理 5.3.14 所需的
弱滤过 $A_\infty$ 双模同态
\[
\varphi\colon
C_1(L^{(1)},L^{(0)};\Lambda_{\mathrm{nov}})
\longrightarrow
C_2(L^{(1)\prime},L^{(0)\prime};
\Lambda_{\mathrm{nov}}).
\]
接下来的两个小节将证明
这一同态确为同伦等价。

须指出，在第 5.3.1 小节，
一旦构造出\emph{滤过}
$A_\infty$ 双模同态，
定理 5.2.35 就立即表明
它是同伦等价。
但引理 5.3.25 所构造的
同态 $\phi$ 仅为\emph{弱滤过}，
因此不能对 $\phi$
应用定理 5.2.35。
为证明 $\phi$ 是同伦等价，
我们必须直接用几何方法
构造同伦逆；
这一构造将在第 5.3.3、
5.3.4 小节完成。
\begingroup\renewcommand{\qedsymbol}{}\end{proof}\endgroup
